\documentclass[11pt]{article}
\usepackage[margin=1in]{geometry}
\usepackage{longtable}
\usepackage{booktabs}

\title{Regional Workforce Dynamics in Texas: A County-Level Assessment}
\author{Policy Research Division}
\date{July 2026}

\begin{document}
\maketitle

\section{Introduction}

Texas continues to add residents faster than any other large state, but the
gains are unevenly distributed across its 254 counties. Metropolitan cores
absorb the bulk of new arrivals while many rural counties face sustained
outmigration of working-age adults. Understanding where labor force
participation is strengthening---and where it is eroding---is essential for
targeting state workforce investments. This brief presents county-level
indicators drawn from our simulated demographic panel and organizes them into
two reference exhibits.

The full county panel appears in Table~\ref{tab:countypanel}, which reports
population, a composite economic vitality index, labor force participation,
and five-year growth for each county in the sample. Because the panel spans
several pages, readers seeking a quick orientation may prefer the compact
subregional summary in Table~\ref{tab:subregions}.

\section{The County Panel}

Three patterns stand out in Table~\ref{tab:countypanel}. First, suburban
collar counties post the strongest participation rates, consistent with their
younger age structures. Second, energy-dependent counties show volatile
vitality scores that track commodity cycles rather than underlying
demographics. Third, a band of counties in the eastern timber region combines
modest population with participation rates well below the state average,
suggesting a mismatch between local job openings and resident skills.

\begin{longtable}{lrrrr}
\caption{Simulated county demographic and workforce panel}\label{tab:countypanel}\\
\toprule
County & Population & Vitality index & LFP (\%) & 5-yr growth (\%) \\
\midrule
\endfirsthead
\caption[]{Simulated county demographic and workforce panel (continued)}\\
\toprule
County & Population & Vitality index & LFP (\%) & 5-yr growth (\%) \\
\midrule
\endhead
\midrule
\multicolumn{5}{r}{\emph{continued on next page}}\\
\endfoot
\bottomrule
\endlastfoot
Nacogdoches & 4,643,670 & 61.8 & 66.8 & 3.9 \\
Bexar & 1,638,156 & 67.3 & 67.2 & 0.5 \\
Travis & 1,971,102 & 65.3 & 44.1 & 7.5 \\
Harris & 4,679,440 & 70.6 & 46.4 & -3.3 \\
Dallas & 4,166,159 & 50.1 & 48.2 & -3.5 \\
Tarrant & 4,458,049 & 55.9 & 65.5 & 7.3 \\
El Paso & 1,594,680 & 62.6 & 71.1 & 1.0 \\
Hidalgo & 2,842,896 & 50.6 & 42.6 & -3.4 \\
Collin & 759,728 & 73.8 & 56.3 & -3.3 \\
Denton & 3,669,271 & 62.4 & 48.9 & 0.0 \\
Fort Bend & 2,287,252 & 76.6 & 49.6 & 7.6 \\
Montgomery & 554,459 & 43.9 & 60.1 & -2.0 \\
Williamson & 1,003,803 & 70.2 & 58.0 & -2.0 \\
Cameron & 770,387 & 92.8 & 54.9 & -1.8 \\
Nueces & 330,041 & 59.3 & 59.9 & 4.2 \\
Brazoria & 248,935 & 78.4 & 61.3 & 4.0 \\
Bell & 2,458,899 & 55.9 & 55.6 & 6.8 \\
Galveston & 314,117 & 90.6 & 64.7 & -1.4 \\
Lubbock & 3,946,353 & 87.7 & 66.1 & -1.9 \\
Webb & 2,967,948 & 88.8 & 59.1 & 5.5 \\
Jefferson & 2,502,732 & 78.0 & 43.3 & 2.6 \\
McLennan & 3,622,442 & 67.0 & 64.7 & -2.2 \\
Smith & 3,808,694 & 59.2 & 57.4 & -0.3 \\
Brazos & 4,432,304 & 84.6 & 55.6 & 6.5 \\
Hays & 3,668,307 & 39.1 & 67.6 & 7.6 \\
Ellis & 3,236,145 & 58.2 & 66.4 & -2.4 \\
Midland & 2,007,370 & 43.9 & 53.4 & 1.7 \\
Ector & 891,704 & 67.9 & 55.2 & -0.5 \\
Johnson & 3,936,679 & 41.7 & 64.3 & 0.7 \\
Guadalupe & 1,386,494 & 45.3 & 59.6 & 1.4 \\
Comal & 571,217 & 93.0 & 62.8 & 3.2 \\
Randall & 1,972,491 & 68.6 & 51.9 & -1.0 \\
Wichita & 1,880,589 & 79.6 & 70.9 & -2.9 \\
Parker & 2,960,260 & 71.1 & 52.2 & -3.8 \\
Grayson & 1,691,932 & 38.3 & 47.0 & -3.0 \\
Gregg & 2,070,534 & 80.5 & 49.9 & 0.2 \\
Taylor & 2,635,174 & 65.5 & 54.9 & 4.4 \\
Potter & 3,195,892 & 49.0 & 45.5 & -0.2 \\
Kaufman & 3,708,788 & 88.3 & 61.6 & 4.8 \\
Rockwall & 892,509 & 59.3 & 57.7 & -1.4 \\
Tom Green & 257,077 & 51.7 & 48.2 & 6.6 \\
Bowie & 4,590,251 & 58.1 & 58.7 & -1.8 \\
Victoria & 3,981,130 & 43.2 & 46.3 & -0.1 \\
Angelina & 4,332,587 & 73.9 & 48.3 & 6.2 \\
Hunt & 171,022 & 91.1 & 62.1 & 2.5 \\
Orange & 3,332,898 & 53.6 & 42.1 & 7.5 \\
Liberty & 2,844,881 & 46.4 & 49.3 & 3.1 \\
Henderson & 1,601,410 & 82.3 & 69.8 & -3.9 \\
Coryell & 1,529,958 & 77.5 & 52.9 & 3.8 \\
Bastrop & 277,291 & 94.1 & 55.9 & -1.4 \\
Walker & 944,914 & 47.8 & 55.0 & 0.1 \\
San Patricio & 4,802,710 & 80.3 & 45.7 & 4.6 \\
Harrison & 1,575,201 & 52.4 & 43.4 & 3.7 \\
Nolan & 815,716 & 38.6 & 52.1 & 6.8 \\
Anderson & 180,141 & 88.0 & 46.1 & 3.6 \\
Wise & 3,726,918 & 53.5 & 65.1 & -2.8 \\
Hardin & 2,777,699 & 46.9 & 45.0 & 6.9 \\
Rusk & 1,647,434 & 91.7 & 63.0 & 8.1 \\
Van Zandt & 3,249,268 & 72.1 & 54.9 & 7.4 \\
Cherokee & 669,991 & 83.0 & 60.3 & -2.9 \\
Kerr & 3,172,903 & 84.8 & 42.5 & -1.0 \\
Lamar & 3,519,590 & 97.2 & 63.2 & 6.4 \\
Maverick & 950,303 & 72.5 & 55.4 & 0.2 \\
Val Verde & 2,255,925 & 54.8 & 48.5 & 3.7 \\
Navarro & 1,227,431 & 54.6 & 63.5 & -2.0 \\
Medina & 2,321,100 & 87.7 & 45.3 & 1.3 \\
Atascosa & 4,808,101 & 49.9 & 62.0 & 0.0 \\
Polk & 2,851,739 & 42.4 & 45.9 & 5.5 \\
Wilson & 3,778,990 & 73.0 & 45.6 & 7.3 \\
Burnet & 2,793,990 & 39.9 & 46.6 & 5.8 \\
Wood & 3,457,628 & 66.6 & 55.9 & 8.1 \\
Starr & 624,517 & 53.2 & 46.9 & 6.4 \\
Wharton & 4,474,834 & 46.8 & 69.5 & -3.5 \\
Erath & 1,809,470 & 61.0 & 71.3 & 0.6 \\
Jim Wells & 4,256,576 & 74.0 & 56.5 & -0.5 \\
Caldwell & 2,177,557 & 71.5 & 44.7 & 4.8 \\
Kendall & 2,439,562 & 42.1 & 54.1 & 3.8 \\
Cooke & 1,925,226 & 46.2 & 46.6 & 1.9 \\
Waller & 390,132 & 48.4 & 44.2 & -0.3 \\
Chambers & 1,878,508 & 97.1 & 66.7 & 7.9 \\
Upshur & 623,159 & 41.3 & 70.4 & 5.4 \\
Howard & 3,404,413 & 62.5 & 63.5 & 6.0 \\
Hopkins & 249,966 & 51.4 & 46.5 & 2.0 \\
Hill & 3,722,290 & 84.9 & 51.9 & 6.1 \\
Fannin & 1,979,197 & 68.0 & 54.2 & -2.6 \\
Hood & 3,451,712 & 82.2 & 53.6 & 1.7 \\
Titus & 909,090 & 93.0 & 59.4 & 6.5 \\
Matagorda & 705,486 & 91.2 & 59.4 & 1.7 \\
Cass & 4,193,469 & 82.9 & 47.4 & 8.6 \\
Grimes & 4,410,682 & 78.3 & 59.1 & 2.9 \\
Uvalde & 1,943,872 & 86.5 & 68.5 & 6.2 \\
Hale & 3,227,254 & 74.3 & 63.3 & 5.9 \\
Bee & 3,159,155 & 82.6 & 66.7 & -0.9 \\
Shelby & 737,845 & 50.2 & 53.4 & 0.2 \\
Fayette & 3,590,599 & 81.8 & 47.7 & 3.9 \\
Aransas & 3,541,369 & 44.7 & 51.8 & -0.8 \\
Palo Pinto & 1,191,616 & 70.4 & 69.2 & -0.7 \\
Panola & 407,781 & 75.6 & 65.6 & 5.0 \\
Milam & 4,524,068 & 43.4 & 61.4 & -2.1 \\
Limestone & 3,847,266 & 47.5 & 44.8 & -2.5 \\
Gillespie & 642,582 & 71.4 & 66.7 & 2.4 \\
Houston & 4,032,751 & 65.1 & 49.5 & -0.9 \\
Austin & 1,025,045 & 55.0 & 69.2 & -1.4 \\
Hutchinson & 2,113,800 & 88.7 & 57.6 & 6.6 \\
Calhoun & 1,458,623 & 82.1 & 69.3 & -3.9 \\
Colorado & 963,144 & 85.2 & 45.3 & 6.2 \\
Gray & 952,062 & 67.1 & 62.9 & -1.8 \\
Tyler & 4,894,257 & 90.9 & 56.8 & 0.9 \\
Moore & 482,513 & 95.3 & 53.5 & 1.6 \\
Willacy & 770,261 & 71.5 & 71.2 & 8.6 \\
Lampasas & 4,305,089 & 63.7 & 62.8 & 6.8 \\
Lavaca & 4,142,182 & 52.3 & 62.9 & 8.2 \\
DeWitt & 1,242,483 & 55.9 & 47.4 & -0.2 \\
Freestone & 2,600,534 & 50.8 & 69.0 & 2.3 \\
Montague & 4,516,772 & 71.6 & 46.1 & 7.8 \\
Gonzales & 4,022,213 & 75.6 & 52.9 & 7.1 \\
Llano & 1,762,763 & 75.0 & 46.1 & 4.6 \\
Jones & 3,544,360 & 52.4 & 45.2 & 6.6 \\
Frio & 3,027,779 & 87.3 & 42.7 & -0.0 \\
Young & 2,500,443 & 70.6 & 52.3 & 2.8 \\
\end{longtable}

\section{Subregional Summary}

Aggregating the panel into eight analytic subregions clarifies the broad
geography of opportunity. Table~\ref{tab:subregions} reports employer counts,
average participation, and the dominant industry cluster for each subregion.
The contrast between the fastest and slowest subregions is stark: the gap in
average participation exceeds twenty percentage points, a difference that no
plausible training intervention could close on its own. Policy responses will
need to pair workforce programs with housing and transportation strategies
that let workers reach the jobs that already exist.

\begin{longtable}{lrrl}
\caption{Subregional workforce summary}\label{tab:subregions}\\
\toprule
Subregion & Employers & Avg.\ LFP (\%) & Dominant cluster \\
\midrule
\endfirsthead
\toprule
Subregion & Employers & Avg.\ LFP (\%) & Dominant cluster \\
\midrule
\endhead
\bottomrule
\endlastfoot
Alamo Corridor & 8,214 & 63.7 & Manufacturing \\
Piney Uplands & 947 & 41.9 & Timber \\
Caprock Rim & 1,533 & 58.2 & Agriculture \\
Trans-Basin & 2,761 & 66.4 & Energy \\
Blackland Arc & 6,089 & 71.8 & Logistics \\
Coastal Bend Sub & 3,412 & 55.1 & Petrochemical \\
Hill Margin & 1,278 & 68.3 & Tourism \\
Red River Shelf & 594 & 49.6 & Ranching \\
\end{longtable}

\section{Conclusion}

The exhibits above are illustrative rather than official statistics, but the
structure of the analysis mirrors what a full census-based panel would
support. A production version of Table~\ref{tab:countypanel} would draw on
the American Community Survey five-year files, and the subregional groupings
in Table~\ref{tab:subregions} would follow the state comptroller's planning
regions. Either way, the policy conclusion is unchanged: participation gaps
are geographic before they are individual, and the state's workforce strategy
should be built accordingly.

\end{document}
